\section{Pedagogy and Course Comments}
\label{sec:pedagogy}

The course provides two different levels of practical database work. In the
earlier mini-project, the team works with a prepared dataset and concentrates
mainly on Entity--Relationship modelling, normalization, implementation and
queries. The present project starts before that stage. The team has to find
suitable data, examine how the information is published, decide which records
can be used and prepare them before database design can be finalized.

This progression makes the larger project useful because many decisions that
appear straightforward with a prepared dataset become less obvious when
working with real source files. Keys, relationships, missing values and even
the meaning of a record often have to be established from the source structure
before they can be represented in the database.

\subsection{Working with Real Data}
\label{subsec:real-data-learning}

Working with real environmental records gives the database concepts taught in
the course a clearer practical meaning. A primary key is not chosen only from
an example relation; it has to match the actual observation represented by the
source. In the same way, a relationship in the ERD is useful only when the
available records contain a dependable way to establish that relationship.

The project also makes the importance of record grain clear. Daily, monthly,
yearly and fiscal-period observations describe different kinds of records and
cannot be treated as interchangeable simply because they belong to the same
environmental sector. Understanding this difference affects extraction,
candidate-key selection, normalization and the final relational structure.

Table~\ref{tab:mini-present-project} shows how the present project extends the
earlier mini-project.

\begin{table}[H]
\centering
\caption{Earlier mini-project and present data-integration project}
\label{tab:mini-present-project}
\setlength{\tabcolsep}{6pt}
\renewcommand{\arraystretch}{1.22}
\begin{tabular}{p{3.2cm}p{4.6cm}p{6.0cm}}
\toprule
\textbf{Aspect} & \textbf{Earlier Mini-Project} & \textbf{Present Project} \\
\midrule
Data & One prepared dataset is provided. & Several real source files have to be
collected, reviewed and selected before implementation. \\
Starting point & Work begins mainly from database modelling. & Work begins with
source discovery, verification and understanding the structure of the
published data. \\
Data preparation & Only limited preparation is required. & Missing values,
irregular layouts, inconsistent names, repeated structures and source-specific
formats have to be handled. \\
Modelling & Entities and relationships are identified from an already prepared
structure. & Candidate entities, keys and relationships are revised as the
actual source records are examined. \\
Normalization & Normalization is applied to a controlled relation. &
Normalization is applied to structures produced from real source data and
observed functional dependencies. \\
ETL & Data preparation before implementation is limited. & Extraction and
transformation form a substantial part of the database development process. \\
Queries & Queries are one of the main implementation exercises. & Queries are
used after source preparation, modelling, normalization, loading and
relationship verification. \\
\bottomrule
\end{tabular}
\end{table}

\subsection{Difficulties Encountered in Practice}
\label{subsec:practical-difficulties}

Most of the difficult work occurs before the final SQLite database becomes
visible. The source files are primarily prepared for human reading rather than
direct database loading. Merged headings, station names outside individual
rows, repeated date columns, inconsistent missing-value notation and several
tables within the same sheet all require interpretation before extraction can
be performed reliably.

Another difficulty is that the different stages of the project depend on one
another. An incorrect interpretation during extraction can produce an
incorrect key; an incorrect key can affect normalization and the ERD; and an
incorrect relationship can later appear as a loading or foreign-key problem.
For this reason, errors cannot always be corrected within only one stage of the
workflow.

The main practical difficulties and corresponding lessons are summarized in
Table~\ref{tab:course-practical-lessons}.

\begin{table}[!htbp]
\centering
\caption{Practical difficulties and lessons from the project}
\label{tab:course-practical-lessons}
\setlength{\tabcolsep}{6pt}
\renewcommand{\arraystretch}{1.22}
\begin{tabular}{p{4.3cm}p{9.5cm}}
\toprule
\textbf{Difficulty} & \textbf{Practical Lesson} \\
\midrule
Irregular source layouts & The structure of a source has to be understood
before extraction is automated. Headings, repeated blocks, dates and station
information cannot always be interpreted from individual rows alone. \\
Inconsistent names and missing values & Standardization should follow
documented rules. A value should not be changed simply to make different files
look consistent. \\
Different observation grains & The unit represented by one record has to be
decided before selecting keys or combining datasets. \\
ERD integration & Similar-looking entities from different sources still have
to be compared by meaning, grain and identifiers before they are merged. \\
Normalization of real records & Functional dependencies become easier to
understand when they are examined against actual extracted records rather than
only theoretical examples. \\
ETL development & Extraction is not simply copying values from one file to
another. Each transformed value has to retain its original meaning and context. \\
Team coordination & Important schema and data decisions need shared review
because a change can affect extraction, normalization, implementation, queries
and documentation at the same time. \\
\bottomrule
\end{tabular}
\end{table}

\subsection{Comments on Course Delivery}
\label{subsec:course-comments}

The progression from a smaller prepared-data exercise to a larger real-data
project works well because students first learn the basic database operations
before dealing with irregular source material. A few additions could make the
transition between the two stages smoother for future groups.

A short demonstration using an actual irregular spreadsheet before the first
sprint would be useful. The demonstration could show how to identify a true
data region, interpret merged headings, recognize repeated columns,
distinguish missing values from zero and determine the grain of a record before
extraction begins. These are small issues individually, but they consume
considerable time when students encounter them for the first time during the
project.

An early review of the first integrated ERD would also be helpful. At that
stage, feedback on candidate keys, record grain, shared entities and proposed
relationships could identify structural problems before they affect
normalization and implementation.

A common review checklist could further improve consistency between groups.
The checklist could cover source authority, record grain, candidate keys,
missing-value treatment, relationship evidence and source traceability. It
would provide a minimum standard for review without requiring different sector
groups to use the same database design. A dedicated final-review period could
then cover database validation, ERD readability, figure quality, source
attribution and unresolved data issues.

Finally, occasional short cross-group review sessions could allow groups to
compare the kinds of source and modelling problems they encounter. Since each
sector contains different data structures, these discussions could expose
students to alternative design decisions while leaving each group responsible
for its own implementation.
